\section{Métricas de evaluación}
\label{sec:metricas}

En el contexto de modelos de \textit{machine learning}, las métricas de evaluación son herramientas fundamentales para cuantificar el desempeño de un modelo en tareas específicas. Estas métricas permiten comparar diferentes modelos, ajustar hiperparámetros y validar la efectividad del aprendizaje. La elección de la métrica adecuada depende del tipo de problema que se esté abordando.


\subsection{Métricas para problemas de regresión}
\label{sec:metricas_regresion}

Algunas métricas comunes para problemas de regresión son:

\begin{itemize}
    \item Error cuadrático medio (MSE, por sus siglas del inglés \textit{Mean Squared Error}): Calcula la media de los cuadrados de las diferencias entre las predicciones y los valores reales. \cite{botchkarev2019performance}
    \begin{equation}
        \text{MSE} = \frac{1}{n} \sum_{i=1}^{n} (y_i - \hat{y}_i)^2
    \end{equation}

    \item Error absoluto medio (MAE, por sus siglas del inglés \textit{Mean Absolute Error}): Calcula la media de las diferencias absolutas entre las predicciones y los valores reales. \cite{willmott2005advantages}
    \begin{equation}
        \text{MAE} = \frac{1}{n} \sum_{i=1}^{n} |y_i - \hat{y}_i|
    \end{equation}

    \item Coeficiente de determinación ($R^2$): Mide la proporción de la varianza en los datos que es explicada por el modelo. Un valor de $R^2$ cercano a 1 indica un buen ajuste. \cite{alexander2016rsquared}
    \begin{equation}
        R^2 = 1 - \frac{\sum_{i=1}^{n} (y_i - \hat{y}_i)^2}{\sum_{i=1}^{n} (y_i - \bar{y})^2}
    \end{equation}
\end{itemize}

\subsection{Métricas para problemas de clasificación}
\label{sec:metricas_clasificacion}
Para el caso de los problemas de clasificación, se suelen utilizar métricas definidas a partir de los siguientes conceptos: \textit{verdaderos positivos} (TP, por sus siglas del inglés \textit{True Positives}), que corresponden al número de instancias positivas correctamente clasificadas como positivas; \textit{falsos positivos} (FP, por sus siglas del inglés \textit{False Positives}), que representan el número de instancias negativas incorrectamente clasificadas como positivas; \textit{verdaderos negativos} (TN, por sus siglas del inglés \textit{True Negatives}), que son el número de instancias negativas correctamente clasificadas como negativas; y \textit{falsos negativos} (FN, por sus siglas del inglés \textit{False Negatives}), que indican el número de instancias positivas incorrectamente clasificadas como negativas. En la \ref{tab:confusion_matrix} se muestra una matriz de confusión \cite{sokolova2009systematic} que resume estos conceptos para un problema de clasificación binaria.

\begin{table}[H]
    \centering
    \caption{Matriz de confusión para un problema de clasificación binaria.}
    \begin{tabular}{c c c}
        \toprule
        & \textbf{Predicción Positiva} & \textbf{Predicción Negativa} \\
        \midrule
        \textbf{Clase Positiva} & TP & FN \\
        \textbf{Clase Negativa} & FP & TN \\
        \bottomrule
        \hline
    \end{tabular}
    \label{tab:confusion_matrix}
\end{table}

\begin{itemize}
    \item Exactitud (\textit{accuracy} en inglés): Mide la proporción de predicciones correctas sobre el total de predicciones. Se computa como se muestra en la \refeq{accuracy} \cite{bishop2006pattern}.
    \begin{equation}
        \label{eq:accuracy}
        \text{accuracy} = \frac{TP + TN}{TP + TN + FP + FN}
    \end{equation}
    \item Precisión (\textit{precision} en inglés): Mide la proporción de predicciones positivas correctas sobre el total de predicciones positivas. Se computa como se muestra en la \refeq{precision} \cite{bishop2006pattern}.
    \begin{equation}
        \label{eq:precision}
        \text{precision} = \frac{TP}{TP + FP}
    \end{equation}
    \item Sensibilidad (\textit{recall} en inglés): Mide la proporción de verdaderos positivos sobre el total de positivos reales. Se computa como se muestra en la \refeq{recall} \cite{bishop2006pattern}.
    \begin{equation}
        \label{eq:recall}
        \text{recall} = \frac{TP}{TP + FN}
    \end{equation}
    \item Puntaje F1 (\textit{F1-score} en inglés): Es la media armónica entre la precisión y la sensibilidad, proporcionando un balance entre ambas métricas. Se computa como se muestra en la \refeq{f1_score} \cite{bishop2006pattern}.
    \begin{equation}
        \label{eq:f1_score}
        \text{F1-score} = 2 \cdot \frac{\text{precision} \cdot \text{recall}}{\text{precision} + \text{recall}}
    \end{equation}
\end{itemize}
